% Editable transcription from the unchanged licensed published PDF.
% Main outcome: Olivier Haution, Theorem 2.1, DOI 10.14231/AG-2015-003.
% Original physical pages 3--6, printed pages 46--49.
% No mathematical proof is generated; source language, order and formulas retained.
\documentclass[11pt]{article}
\usepackage[a4paper,margin=25mm]{geometry}
\usepackage{fontspec}
\setmainfont{lmroman10-regular.otf}[BoldFont=lmroman10-bold.otf,ItalicFont=lmroman10-italic.otf,BoldItalicFont=lmroman10-bolditalic.otf]
\usepackage{amsmath,amssymb,amsthm,amscd}
\usepackage{xurl,hyperref}
\hypersetup{hidelinks}
\newcommand{\OO}{\mathcal O}
\newcommand{\Sing}{\operatorname{Sing}}
\newcommand{\codim}{\operatorname{codim}}
\newcommand{\Sch}{\operatorname{Sch}}
\newcommand{\Spec}{\operatorname{Spec}}
\newcommand{\Proj}{\operatorname{Proj}}
\newcommand{\trdeg}{\operatorname{tr.deg.}}
\newtheorem*{mainthm}{Theorem 2.1}
\newtheorem*{supportcor}{Corollary 2.3}
\newtheorem*{supportlem}{Lemma 2.4}
\newtheorem*{sourceremark}{Remark 2.2}
\setlength{\emergencystretch}{1em}
\begin{document}
\begin{center}
{\Large 3215: Resolving singularities in codimension two}\par\medskip
Olivier Haution
\end{center}
\noindent\textbf{Problem.} Prove Theorem 2.1 below, with the original conventions and human-authored proof transcribed in full.

\smallskip
\noindent\textbf{Source and attribution.} Olivier Haution, \emph{Detection by regular schemes in degree two}, Algebraic Geometry \textbf{2} (1) (2015), 44--61, DOI \href{https://doi.org/10.14231/AG-2015-003}{10.14231/AG-2015-003}. This editable transcription uses the unchanged official published version, physical PDF pages 3--6 (printed 46--49), plus the cited bibliography entries from physical 17 (printed 60). Copyright Foundation Compositio Mathematica 2015. Reused under the \href{https://creativecommons.org/licenses/by-nc/3.0/}{Creative Commons Attribution-NonCommercial 3.0 Unported license}: attribution and noncommercial reuse are required. Mathematical prose and formulas are transcribed; typesetting, page breaks, heading style and citation links are changed. The selected outcome is Theorem 2.1 in full. No complete resolution in every codimension is claimed. Corollary 2.3 and Lemma 2.4 below are supporting material, not additional problems.

\noindent\textbf{Editorial difficulty note (not author text).} Difficulty H2. The proof must globalize local two-dimensional resolutions while controlling every point lying over codimension at most two. It constructs a global blowup center through scheme-theoretic images and verifies its local models using flat base change and regularity descent. Passing from normal to reduced quasi-excellent schemes additionally requires finite-normalization and codimension control, rather than a direct invocation of local resolution.

% SOURCE BLOCK C1: original physical 3, printed 46, sections 1.1--1.6.
\section*{1. Terminology and conventions}
\subsection*{1.1 Schemes and morphisms}
All schemes are noetherian and separated. When $S$ is a scheme, we denote by $\Sch_S$ the category of schemes of finite type over $S$ and proper $S$-morphisms.

We say that a morphism $Y\to X$ is, or induces, an isomorphism outside of a closed subset $Z$ of $X$ if the base change $(X-Z)\times_XY\to X-Z$ is an isomorphism.

\subsection*{1.2 Regularity}
We say that a scheme $X$ is \emph{regular at a point $x$} when the local ring $\OO_{X,x}$ is regular. Otherwise we say that $x$ is a \emph{singular point} of $X$. The \emph{singular locus} $\Sing X$ is the set of such points. A scheme is \emph{regular} if it is regular at each of its points.

\subsection*{1.3 Normalisation}
When $X$ is an integral scheme, its \emph{normalisation} $X'\to X$ is the affine morphism given locally by the integral closure of its coordinate ring in its function field. More generally, when $X$ is an arbitrary scheme, its normalisation is the composition $\coprod_iX_i'\to\coprod_iX_i\to X$, where the $X_i$ are the irreducible components of $X$, and $X_i'\to X_i$ their normalisations.

\subsection*{1.4 Excellence}
A scheme $S$ is called \emph{quasi-excellent} if
\begin{itemize}
\item any integral scheme finite over $S$ admits an open dense regular subscheme,
\item and for any closed point $x$ of $X$, the completion morphism of $\Spec\OO_{X,x}$ has geometrically regular fibres.
\end{itemize}
If $S$ is quasi-excellent, then any scheme of finite type over $S$ is quasi-excellent, and so is the localisation of $S$ at any of its points [ILO12, \S2, Theorem 5.1]. The spectrum of a field, and of $\mathbb Z$, is quasi-excellent, so that the class of quasi-excellent schemes contains all varieties and arithmetic schemes. The normalisation of a quasi-excellent scheme is finite [EGA IV, (7.6.1)]; the singular locus of a quasi-excellent scheme is closed [EGA IV, (6.12.3)]. For a scheme of finite type over a regular scheme, we will use the terminology \emph{excellent} instead of quasi-excellent; this is compatible with [EGA IV, (7.8.2)].

\subsection*{1.5 Projectivity}
A morphism $Y\to X$ is \emph{projective} if there are a quasi-coherent sheaf of graded $\OO_X$-algebras $S_\bullet$, with $S_1$ a coherent $\OO_X$-module generating $S_\bullet$ as an $\OO_X$-algebra, and an isomorphism over $X$ between $Y$ and $\Proj_XS_\bullet$. A morphism is \emph{quasi-projective} if it decomposes as an open immersion followed by a projective morphism.

\subsection*{1.6 Envelopes}
A proper morphism $f:Y\to X$ is an \emph{envelope} if for any integral closed subscheme $Z$ of $X$, there is a closed subscheme $W$ of $Y$ such that $f$ induces a birational morphism $W\to Z$. It is equivalent to require that $Y(K)\to X(K)$ be surjective for all fields $K$. When $X$ is of finite type over $S$, a \emph{Chow envelope} over $S$ is an envelope $Y\to X$ with $Y$ quasi-projective over $S$. By Chow's Lemma [EGA II, (5.6.1)] and noetherian induction, such an envelope always exists.

% SOURCE BLOCK C2: original physical 4, printed 47, full section 1.7.
\subsection*{1.7 Dimension}
The codimension $\codim_X(x)$ of a point $x$ in a scheme $X$ is the dimension of the local ring $\OO_{X,x}$. When $Y$ is a closed subset of $X$, its codimension $\codim_X(Y)$ is the infimum of the codimensions in $X$ of the points of $Y$. When $X$ is an integral scheme of finite type over a scheme $S$, the dimension of $X$ over $S$ is defined as
\[
\dim_SX=\trdeg(\kappa(X)/\kappa(\overline X))-\codim_S(\overline X),
\]
where $\overline X$ is the scheme-theoretic image of $X\to S$, and $\kappa(X)$ and $\kappa(\overline X)$ are the function fields of $X$ and $\overline X$, respectively. When $Z$ is an integral closed subscheme of $X$, we have, using the dimension formula [EGA IV, (5.6.5.1)],
\begin{align*}
\dim_SZ&=\trdeg(\kappa(Z)/\kappa(\overline Z))-\codim_S(\overline Z)\\
&\leq\trdeg(\kappa(Z)/\kappa(\overline Z))-\codim_X(Z)-\codim_S(\overline X)\\
&\leq\trdeg(\kappa(X)/\kappa(\overline X))-\codim_X(Z)-\codim_S(\overline X)\\
&=\dim_SX-\codim_X(Z).
\end{align*}
In particular, $\dim_SZ\leq\dim_SX$, so that when $Y$ is a scheme of finite type over $S$, we may define $\dim_SY$ as the supremum of the integers $\dim_SZ$ for $Z$ an integral closed subscheme of $Y$.

Let $X$ be an integral scheme of finite type over $S$ and let $Y$ be a closed subscheme of $X$. Since for any integral closed subscheme $Z$ of $Y$, we have $\dim_SZ\leq\dim_SX-\codim_X(Z)\leq\dim_SX-\codim_X(Y)$, it follows that
\begin{equation}
\dim_SY\leq\dim_SX-\codim_X(Y).\tag{1}
\end{equation}

% SOURCE BLOCK T1: original physical 4--5, printed 47--48, full Theorem 2.1.
\section*{2. Resolving singularities in codimension two}
\begin{mainthm}
Let $X$ be a reduced quasi-excellent scheme. Then there are a reduced scheme $\widetilde X$ and a projective birational morphism $f:\widetilde X\to X$ with the following properties.
\begin{enumerate}
\item[(i)] The scheme $\widetilde X$ is regular at any of its points $x$ such that $f(x)$ has codimension at most two in $X$.
\item[(ii)] The morphism $f$ is an isomorphism outside of the singular locus of $X$.
\end{enumerate}
\end{mainthm}

% SOURCE BLOCK P1: original physical 5--6, printed 48--49, full proof of 2.1.
\begin{proof}
We first assume that $X$ is normal. Then $\Sing X$ contains no point of codimension at most one in $X$. Let $\Sigma$ be the set of points of $\Sing X$ which have codimension two in $X$. We view $\Sigma$ and $\Sing X$ as topological subspaces of $X$. Then $\Sigma$ consists of generic points of $\Sing X$. Since $X$ is quasi-excellent, $\Sing X$ is the support of a noetherian scheme, and therefore the space $\Sigma$ is finite and discrete. Thus we can write $\Sigma=\{s_1,\ldots,s_n\}$, with $s_i$ not belonging to the closure of $\{s_j\mid j\ne i\}$ in $X$.

For every $i\in\{1,\ldots,n\}$, the scheme $Y_i=\Spec\OO_{X,s_i}$ is a singular normal quasi-excellent scheme of dimension two. By [Lip78, Theorem p.~151, and C p.~155], there is a zero-dimensional closed subscheme $Z_i$ of $Y_i$ such that the blow-up $g_i:\widetilde Y_i\to Y_i$ of $Y_i$ along $Z_i$ is a desingularisation; that is, $\widetilde Y_i$ is regular.

We let $Z=\coprod_{i=1}^nZ_i$ and $Y=\coprod_{i=1}^nY_i$, and claim that $z:Z\to Z\times_XY$ is an isomorphism. Indeed, since the open subschemes $Z_j\times_XY_i$ for $i,j=1,\ldots,n$ cover $Z\times_XY$, it suffices to see that the base change $z_{i,j}:Z_j\times_YY_i\to Z_j\times_XY_i$ of $z$ is an isomorphism. If $j\ne i$, no point of $Y_i$ is mapped to $s_j$, so that $Z_i\times_XY_j=\varnothing$ and $z_{i,j}$ is an isomorphism. The morphism $Y_i\to X$ is a monomorphism by [EGA I, I (2.4.2)], and so is $Y_i\to Y$. Therefore $d_i:Z_i\to Z_i\times_YY_i$ and $z_{i,i}\circ d_i:Z_i\to Z_i\times_XY_i$ are isomorphisms, and we deduce that $z_{i,i}$ is an isomorphism.

Let $W$ be the scheme-theoretic image of $Z\to X$ and let $f:\widetilde X\to X$ be the blow-up of $X$ along $W$. The morphism $Y\to X$ is flat, and by the compatibility of scheme-theoretic images with flat base change [EGA IV, (2.3.2)], the scheme-theoretic image of the morphism $Z=Y\times_XZ\to Y$ is $Y\times_XW$. Since this morphism is a closed embedding, it follows that $Z\to Y\times_XW$ is an isomorphism. By the compatibility of blow-ups with flat base change, we see that the base change of $f$ along $Y\to X$ is the blow-up of $Y$ along $Z$. In particular, we have a cartesian square, for every $i\in\{1,\ldots,n\}$,
\begin{equation}
\begin{CD}
\widetilde Y_i @>{g_i}>> Y_i\\
@V{u_i}VV @VV{l_i}V\\
\widetilde X @>{f}>> X
\end{CD}\tag{2}
\end{equation}

The set-theoretic image of $Z\to X$ is $\Sigma$ (this is where we use that each $Z_i$ is zero-dimensional), hence $W$ is supported on the closure of $\Sigma$, which is contained in $\Sing X$. Therefore $f$ is an isomorphism outside of $\Sing X$, and in particular is birational. It is projective, as is any blow-up. Moreover, $\widetilde X$ is reduced, being the blow-up of a reduced scheme.

Thus, to conclude the proof in the normal case, it will suffice to verify condition (i). So let $x$ be a point of $\widetilde X$ such that $f(x)$ has codimension at most two in $X$. If $f(x)\notin\Sigma$, then $f(x)\notin\Sing X$ by construction of $\Sigma$. Therefore $f$ induces an isomorphism $\OO_{X,f(x)}\to\OO_{\widetilde X,x}$, and it follows that $\widetilde X$ is regular at $x$. Now assume that $f(x)=s_i$ for some $i\in\{1,\ldots,n\}$. Let $F$ be the residue field of $\widetilde X$ at $x$. Composing the $F$-point given by $x$ with the morphism $f$, we obtain an $F$-point of $X$. Since its image is the point $s_i$, this $F$-point factors through the localisation $l_i:Y_i\to X$. In view of the cartesian square (2), we obtain an $F$-point $y$ of $\widetilde Y_i$ such that $x=u_i(y)$. Since $\widetilde Y_i$ is regular at $y$, and $u_i$ is flat, it follows from [EGA IV, (6.5.2, (i))] that $\widetilde X$ is regular at $x$.

Now let $X$ be an arbitrary reduced quasi-excellent scheme. Its normalisation $\pi:X'\to X$ is finite and birational. Applying the construction above to the quasi-excellent normal scheme $X'$, we obtain a projective birational morphism $f':\widetilde X\to X'$ satisfying conditions (i) and (ii). The morphism $f=\pi\circ f'$ is projective and birational. Since $\pi$ induces an isomorphism outside of $\Sing X$, and $f'$ induces an isomorphism outside of $\Sing X'\subset\pi^{-1}\Sing X$, we see that $f$ satisfies condition (ii). To see that $f$ satisfies condition (i), let $x$ be a point of $\widetilde X$ such that $f(x)$ has codimension at most two in $X$. By Lemma 2.4 below applied to the morphism $\pi$, the point $f'(x)$ has codimension at most two in $X'$. Since $f'$ satisfies (i), it follows that $\widetilde X$ is regular at $x$.
\end{proof}

% SOURCE BLOCK S1: original physical 6, printed 49, Remark 2.2 and full Corollary 2.3.
\begin{sourceremark}
From the proof, we see that we may take for $f$ the normalisation of $X$, followed by the blow-up along a closed subscheme of codimension two.
\end{sourceremark}

We will be using the following variant of Theorem 2.1.

\begin{supportcor}
Let $X$ be a reduced quasi-excellent scheme, of finite type over a scheme $S$. Then there is a projective birational morphism $f:X'\to X$, with $X'$ reduced and quasi-projective over $S$, and regular at any of his points $x$ such that $f(x)$ has codimension at most two in $X$.
\end{supportcor}
\begin{proof}
Applying Chow's lemma [EGA II, (5.6.1)], we obtain a projective birational morphism $Y\to X$ with $Y$ reduced and quasi-projective over $S$. By Lemma 2.4, a point has codimension at most two in $Y$ as soon as its image has codimension at most two in $X$. The statement follows by applying Theorem 2.1 to the scheme $Y$.
\end{proof}

% SOURCE BLOCK S2: original physical 6, printed 49, full Lemma 2.4 and proof.
\begin{supportlem}
Let $f:Y\to X$ be a birational morphism of finite type, and let $y$ be a point of $Y$ with image $x=f(y)$ in $X$. Then we have $\codim_Y(y)\leq\codim_X(x)$.
\end{supportlem}
\begin{proof}
\sloppy
Let $Y_0$ be an irreducible component of $Y$ containing $y$ and such that $\codim_Y(y)=\allowbreak\codim_{Y_0}(y)$. There is a unique irreducible component $X_0$ of $X$ such that $f$ induces a birational morphism $f_0:Y_0\to X_0$ of finite type. Applying the dimension formula [EGA IV, (5.6.5.1)] to the morphism $f_0$, we obtain $\codim_{Y_0}(y)\leq\codim_{X_0}(x)$, and we conclude since $\codim_{X_0}(x)\leq\codim_X(x)$.
\end{proof}

% SOURCE BLOCK R1: original physical 17, printed 60, only entries cited in these blocks.
\section*{References cited in the transcribed text}
\begin{description}
\item[EGA I] A. Grothendieck, \emph{Éléments de géométrie algébrique I, Le langage des schémas}, Publ. Math. Inst. Hautes Études Sci. 4 (1960).
\item[EGA II] A. Grothendieck, \emph{Éléments de géométrie algébrique II, Étude globale élémentaire de quelques classes de morphismes}, Publ. Math. Inst. Hautes Études Sci. 8 (1961).
\item[EGA IV] A. Grothendieck, \emph{Éléments de géométrie algébrique IV, Étude locale des schémas et des morphismes de schémas}, Publ. Math. Inst. Hautes Études Sci. 20 (1964), 24 (1965), 28 (1966) and 32 (1967).
\item[ILO12] L. Illusie, Y. Laszlo, and F. Orgogozo, \emph{Travaux de Gabber sur l'uniformisa\-tion locale et la cohomologie étale des schémas quasi-excellents. Séminaire à l'École polytechnique 2006--2008}, with the collaboration of F. Déglise, A. Moreau, V. Pilloni, M. Raynaud, J. Riou, B. Stroh, and M. Temkin, Astérisque 363--364, Soc. Math. de France, Paris, 2014.
\item[Lip78] J. Lipman, \emph{Desingularization of two-dimensional schemes}, Ann. Math. (2) 107 (1978), no. 1, 151--207.
\end{description}

\section*{Transcription notes}
\noindent\textbf{Difficulty: H2.} The selected proof globalizes two-dimensional local resolutions through a single scheme-theoretic blowup centre, uses flat base change to recover each local resolution, and tracks regularity and codimension through finite normalization. This is above an ordinary PhD qualifying-examination exercise.

\noindent This section is editorial and is separate from the authored mathematics. Original section 1.4 prints ``closed point $x$ of $X$'' after introducing $S$; that source wording is retained. In the cross-localisation paragraph, the source first uses $Z_j\times_XY_i$ and later writes $Z_i\times_XY_j=\varnothing$ for distinct indices; both formulas are retained. No formula, argument or index has been repaired. The selected theorem has no omitted proof step in this source. Lipman's local desingularisation result and the cited EGA and ILO foundations remain referenced external prerequisites. The unrelated source sections 1.8--1.9 and 3--6, and uncited bibliography entries, are outside this transcription.
\end{document}
